\hypertarget{discussion-and-outlook}{%
\section{Discussion and outlook}\label{discussion-and-outlook}}

\hypertarget{what-the-corrections-change}{%
\subsection{What the corrections change}\label{what-the-corrections-change}}

The strongest lesson of the verification campaign is methodological. A finite computation can be exact and still answer
the wrong typed question: aliases can corrupt a denominator; a branch witness can be mistaken for a structure property;
a distance can be called a commutator; a null direction can be quotient gauge; or a solver budget can masquerade as an
existence boundary. Version 2 retains results only after those distinctions are made explicit.

The positive inventory remains useful. The repaired gauge carrier yields a branch-level selection under declared caps;
the regular \(SU(5)\) frame and its coset are explicit exact objects; the content quotient produces one selected orbit;
the \(N\)-copy probe computes real generic ranks; the Boolean access lattice supplies an exact incomparability theorem;
and the min-cut carrier supplies an exact LP dual certificate plus independent contracted-state evidence. None alone is
a theory of nature, but each is a reproducible construction with a clear failure condition.

\hypertarget{nearest-upgrades-precisely-typed-open-problems}{%
\subsection{Nearest upgrades}\label{nearest-upgrades-precisely-typed-open-problems}}

Three upgrades now dominate. First, the 13 residual P1 carriers need either a further exact reduction beyond the five named classes
or a completeness theorem for an exact quotient class; independently, a contracted-state underdetermination
example must be built if the question is to move beyond cuts. Second, BMV relevance requires a dynamical model deriving
record formation, rather than stipulating a measure-and-record channel. Third, the record-stability program requires the
physical invariant algebra and dynamics: UV-to-residual restriction ranks, Grassmann/EOM/syzygy relations, and a
persistence criterion for candidate records.

Other bounded upgrades remain worthwhile: repair the step43 final-carrier counts; replace F-008 summary-reading
validators; establish or refute large-bond RT saturation with a rate; widen the representation and scalar windows; and
generate a physically motivated class of measurement models rather than a two-exemplar taxonomy. A third substrate
could test whether the vocabulary remains useful, but would not by itself prove universality.

\hypertarget{what-kind-of-theory-is-this}{%
\subsection{What kind of framework is this?}\label{what-kind-of-theory-is-this}}

The calculus is not dynamical: it has no Lagrangian, cross-sections, or continuum limit. Its present evidential role is a
finite structural language for asking whether one readout factors through another, whether a candidate set closes, and
which assumptions a recognition imports. The campaign shows both the value and the limit of that language. It can make
hidden quotients and quantifiers visible; it cannot turn an unbuilt physical bridge into a prediction.

\hypertarget{an-invitation}{%
\subsection{An invitation}\label{an-invitation}}

The most useful independent attacks are now concrete: regenerate the repaired carrier under a different declared
quotient; search the 13 P1 residuals for a sixth exact reduction; construct state-level equal-entanglement examples after
physical quotienting; derive a gravitational record channel from dynamics; or compute the restriction and relation
structure of the record invariant ring. A disagreement with any exported finite table is a direct reproducibility
failure. A physical result addresses the open program only when the missing carrier-to-physics map is supplied.
